\documentclass[11pt]{article}
\usepackage{amsmath,amsthm,amssymb}
\usepackage[margin=1in]{geometry}
\usepackage{enumitem}
\usepackage{booktabs}

\newtheorem{theorem}{Theorem}
\newtheorem{lemma}[theorem]{Lemma}
\newtheorem{proposition}[theorem]{Proposition}
\newtheorem{corollary}[theorem]{Corollary}
\newtheorem{conjecture}[theorem]{Conjecture}
\theoremstyle{definition}
\newtheorem{definition}[theorem]{Definition}
\newtheorem{remark}[theorem]{Remark}
\newtheorem{example}[theorem]{Example}

\title{The rank-corner formula at $r=2$:\\
  $\dim B_{\ell+1}^{(\lambda)} V_\lambda = 2$ for $\lambda = (3, 2, 1^{n-5})$}
\author{Clio}
\date{12 May 2026}

\begin{document}
\maketitle

\begin{abstract}
We prove that $\dim B_{n-2}^{(\lambda)} V_\lambda = 2$ for every
$n \ge 8$ at $\lambda = (3, 2, 1^{n-5})$. This closes the dimension
gap left open in the May-20 paper \texttt{2026-05-20-3-2-1n-family.tex},
which proved $E_1^-$ containment (hence rank-zero) but only the upper
bound $\dim \le 2$. The lower bound is the
\textbf{first proven instance of the rank-corner formula at $r=2$
length-preserving corners}.

The argument is a \emph{position-of-$n$ invariant}: the operators
$T_j$ for $j \le n-2$ preserve the position of the letter $n$ in every
SYT. So if $\pi^{(c)}: V_\lambda \to V_\lambda$ is the seminormal-basis
projection onto SYTs with $n$ at the cell $c$, then $\pi^{(c)}$
commutes with each $(T_j{+}1)$ for $j \le n-2$. Under the May-20
three-branch reduction $B_{n-2}^{(\lambda)} V_\lambda =
\langle F_A \rangle + \langle F_B \rangle$, the $A$-piece $F_A$ has
support only on SYTs with $n \in \{(1,3),(n-3,1)\}$, and the
$B$-piece $F_B$ has support on $\{(2,2),(n-3,1)\}$. The cells
$(1,3)$ and $(2,2)$ are exclusive, so it suffices to show
$\pi^{(1,3)} F_A \ne 0$ and $\pi^{(2,2)} F_B \ne 0$. A four-line
Hoefsmit case analysis on the $\mu_X$-supports does each.

\textbf{Verification.} $\dim B_{n-2}^{(\lambda)} V_\lambda = 2$
verified at $q=7$ for $n \in \{8, 9, 10\}$; the exclusive
$\pi^{(1,3)}$-witnesses (4 SYTs) and $\pi^{(2,2)}$-witnesses (5 SYTs)
match the case-analysis prediction exactly.
\end{abstract}

\section{Setup and statement}\label{sec:setup}

Notation as in \texttt{2026-05-20-3-2-1n-family.tex} (henceforth
\textbf{May-20}). Throughout, $n \ge 8$ and
\[
   \lambda := (3, 2, 1^{n-5}) \vdash n,
   \qquad \ell := \ell(\lambda) = n-3,
   \qquad j_0 := \ell+1 = n-2.
\]
The three corners of $\lambda$ are $c_1 = (1, 3)$, $c_2 = (2, 2)$
(both length-preserving), and $c_3 = (n-3, 1)$ (column-1 bottom).
For each pair $\{c_i, c_j\}$, the doubly-stripped partition is
\[
   \mu_A := (2, 2, 1^{n-6}) \vdash n-2,\qquad
   \mu_B := (3, 1^{n-5}) \vdash n-2,\qquad
   \mu_C := (2, 1^{n-4}) \vdash n-2,
\]
indexed so $\mu_A$ omits $\{c_1, c_3\}$, $\mu_B$ omits
$\{c_2, c_3\}$, and $\mu_C$ omits $\{c_1, c_2\}$.

We rely on three external results:

\begin{theorem}[May-20 three-branch reduction, Theorem~3.5]\label{thm:reduction}
\[
   B_{n-2}^{(\lambda)} V_\lambda
   = (T_{n-3}{+}1)(T_{n-2}{+}1)\,
     \big(\Phi_A(u^{(\mu_A)}) + \Phi_B(u^{(\mu_B)})\big),
\]
where $u^{(\mu_X)}$ spans $B_{n-3}^{(\mu_X)} V_{\mu_X}$ and the
$\mu_C$-piece vanishes via May-20 Lemma~4.1.
\end{theorem}

\begin{theorem}[Support of $u^{(\mu_A)}$; May-13 Theorem~2.5 +
May-18]\label{thm:uA}
$u^{(\mu_A)}$ has seminormal support exactly
\[
   S_A := \{(n{-}5, n{-}4),\ (n{-}5, n{-}3),\ (n{-}4, n{-}2),\
   (n{-}3, n{-}2)\}
\]
in the $(\alpha, \gamma)$-indexing of May-13 (column-$2$ entries of
$\mu_A$), with all four coefficients in $\mathbb{Q}(q)^\times$.
\end{theorem}

\begin{theorem}[Support of $u^{(\mu_B)}$; May-12 Proposition~6.5]\label{thm:uB}
$u^{(\mu_B)}$ has seminormal support exactly
\[
   S_B := \{(n{-}5, n{-}4),\ (n{-}5, n{-}3),\ (n{-}4, n{-}3),\
   (n{-}4, n{-}2),\ (n{-}3, n{-}2)\}
\]
in the $(a, b)$-indexing of May-12 (row-$1$ entries at $(1,2),(1,3)$
of $\mu_B$), with all five coefficients in $\mathbb{Q}(q)^\times$.
\end{theorem}

The main result of this paper:

\begin{theorem}[Rank-corner formula at $r=2$]\label{thm:main}
For every $n \ge 8$, $\dim B_{n-2}^{(\lambda)} V_\lambda = 2$.
Equivalently, the rank of $B_{\ell+1}^{(\lambda)}$ on $V_\lambda$
equals the number of length-preserving corners of $\lambda$.
\end{theorem}

May-20 Theorem~6.4 gives $\dim \le 2$; this paper proves
$\dim \ge 2$.

\section{The position-of-$n$ invariant}\label{sec:invariant}

\begin{lemma}[Position of $n$ is invariant under $T_j$ for $j \le n-2$]\label{lem:n-invariant}
For every SYT $T$ of $\lambda$ and every $1 \le j \le n-2$, the
position $\mathrm{pos}_T(n)$ of the letter $n$ is unchanged under
the action of $T_j$ on $v_T$. More precisely, $T_j v_T$ is a
$\mathbb{Q}(q)$-linear combination of seminormal basis vectors
$v_{T'}$ with $\mathrm{pos}_{T'}(n) = \mathrm{pos}_T(n)$.
\end{lemma}

\begin{proof}
$T_j$ acts in the seminormal basis by the standard formula: same
row $\Rightarrow q$, same column $\Rightarrow -1$, otherwise a
Hoefsmit $2$-block on $\{v_T, v_{s_j T}\}$ where $s_j T$ swaps the
letters $j, j+1$. In all three cases, $T_j v_T$ is supported on
SYTs related to $T$ by either the identity or the swap of $j$ and
$j+1$. Since $j+1 \le n-1$, neither swap touches the letter $n$,
so $\mathrm{pos}(n)$ is preserved.
\end{proof}

\begin{corollary}[Position-$n$ projections commute with $(T_j{+}1)$]\label{cor:proj-commute}
For each cell $c$ of $\lambda$, define the seminormal-basis
projection
\[
   \pi^{(c)}: V_\lambda \to V_\lambda,
   \qquad
   \pi^{(c)} v_T = \begin{cases} v_T & \mathrm{pos}_T(n) = c \\ 0 & \text{otherwise.}
   \end{cases}
\]
Then $\pi^{(c)}$ commutes with $(T_j{+}1)$ for every $1 \le j \le n-2$.
\end{corollary}

\begin{proof}
Immediate from Lemma~\ref{lem:n-invariant}, since $T_j$ preserves
each summand $V_\lambda = \bigoplus_c \mathrm{im}(\pi^{(c)})$, and
hence so does $T_j + 1$.
\end{proof}

\section{Support of $F_A$ and $F_B$ under $\pi^{(c)}$}\label{sec:supports}

Define
\[
   F_A := (T_{n-3}{+}1)(T_{n-2}{+}1)\,\Phi_A(u^{(\mu_A)}),
   \qquad
   F_B := (T_{n-3}{+}1)(T_{n-2}{+}1)\,\Phi_B(u^{(\mu_B)}).
\]
By Theorem~\ref{thm:reduction}, $B_{n-2}^{(\lambda)} V_\lambda = \mathbb{Q}(q) F_A + \mathbb{Q}(q) F_B$.

\begin{lemma}[$n$-position support of $F_A$ and $F_B$]\label{lem:supp-FA-FB}
Every SYT $T \in \mathrm{supp}(F_A)$ has $\mathrm{pos}_T(n)
\in \{c_1, c_3\} = \{(1,3), (n-3,1)\}$.
Every SYT $T \in \mathrm{supp}(F_B)$ has $\mathrm{pos}_T(n)
\in \{c_2, c_3\} = \{(2,2), (n-3,1)\}$.
\end{lemma}

\begin{proof}
By construction of $\Phi_A$ (May-20 Definition~2.2), every SYT in
the seminormal support of $\Phi_A(V_{\mu_A})$ has $\{n-1, n\}$
occupying the cell pair $\{c_1, c_3\}$ in one of the two orderings.
In particular, $\mathrm{pos}(n) \in \{c_1, c_3\}$ for every such SYT.
By Corollary~\ref{cor:proj-commute}, this is preserved under
$(T_{n-3}{+}1)(T_{n-2}{+}1)$. The same argument applies to $\Phi_B$
with $\{c_2, c_3\}$.
\end{proof}

\begin{corollary}[Cross-projection vanishings]\label{cor:cross}
$\pi^{(c_1)} F_B = 0$ and $\pi^{(c_2)} F_A = 0$.
\end{corollary}

\begin{proof}
$F_B$'s support has $\mathrm{pos}(n) \in \{c_2, c_3\}$, never $c_1$.
Similarly $F_A$'s support never has $\mathrm{pos}(n) = c_2$.
\end{proof}

\section{$\pi^{(c_1)} F_A \ne 0$ (Lemma A)}\label{sec:LemmaA}

In this section we identify $V_\lambda|_{\mathrm{pos}(n) = c_1}$
with the seminormal module $V_{\lambda^-}$ for
$\lambda^- := \lambda \setminus c_1 = (2, 2, 1^{n-5}) \vdash n-1$ as
$H_q(S_{n-1})$-modules: each SYT $T$ of $\lambda$ with $n$ at
$(1, 3)$ corresponds to the SYT $T'$ of $\lambda^-$ on letters
$1, \ldots, n-1$ obtained by deleting the cell $(1, 3)$.

For each $(\alpha, \gamma) \in S_A$, write
$T^{A,B}_{(\alpha, \gamma)}$ for the SYT of $\lambda$ obtained from
the $\mu_A$-SYT with column-$2$ entries $(\alpha, \gamma)$ by
$\Phi_A$ in the ``$n$ at $c_1$'' direction --- i.e.,
$n$ at $(1, 3)$, $n-1$ at $(n-3, 1)$.

We use the Hoefsmit $2$-block in the form: for $T_i$ acting
on $\{v_T, v_{s_iT}\}$ at axial distance $d = c(i{+}1) - c(i)$ with
$|d| \ge 2$, write
\begin{equation}\label{eq:hoefsmit}
   T_i v_T = a_d v_T + b_d v_{s_iT},
   \qquad
   (T_i{+}1) v_T = (a_d{+}1) v_T + b_d v_{s_iT}.
\end{equation}
The eigenvalues of $T_i$ on this 2-block are $\{q, -1\}$, so the
trace gives $a_d + a'_d = q-1$ (where $a'_d$ is the diagonal entry
on $v_{s_iT}$) and the determinant gives $a_d a'_d - b_d b'_d = -q$.
\emph{Non-degeneracy} of the 2-block (i.e., $b_d \ne 0$ and
$b'_d \ne 0$) holds iff $|d| \ne 1$
(see May-13 \S 2.2 / Murphy's seminormal basis); under
non-degeneracy, $a_d \ne -1$ as well, since $a_d = -1$ would force
$a'_d = q$ and then $a_d a'_d = -q$, requiring $b_d b'_d = 0$, a
contradiction. Hence for $|d| \ge 2$ both $a_d + 1 \in \mathbb{Q}(q)^\times$
and $b_d \in \mathbb{Q}(q)^\times$.

\subsection{Case analysis on $S_A$}

We decompose $\pi^{(c_1)} F_A$ as a sum over the four supports of
$u^{(\mu_A)}$. Since $\Phi_A(v_{T'}^{(\mu_A)}) =
v_{T^{A,A}_{(\alpha,\gamma)}} + \tau_A v_{T^{A,B}_{(\alpha,\gamma)}}$
with $\tau_A \in \mathbb{Q}(q)^\times$ (May-20 Definition~2.2, axial
distance $|d_A| = n-2 \ge 6$ at $n \ge 8$), the $c_1$-projection
selects the $\tau_A$ piece:
\[
   \pi^{(c_1)} \Phi_A(u^{(\mu_A)})
   = \tau_A \sum_{(\alpha, \gamma) \in S_A}
       A_{(\alpha, \gamma)} v_{T^{A,B}_{(\alpha, \gamma)}},
\]
where $A_{(\alpha, \gamma)} \in \mathbb{Q}(q)^\times$ are the
coefficients of $u^{(\mu_A)}$ (Theorem~\ref{thm:uA}).

For each $(\alpha, \gamma) \in S_A$ we compute
$(T_{n-3}{+}1)(T_{n-2}{+}1) v_{T^{A,B}_{(\alpha, \gamma)}}$.
The positions of letters $n-3, n-2$ in
$T^{A,B}_{(\alpha, \gamma)}$ depend on the $\mu_A$-data:

\paragraph{Positions in $T^{A,B}_{(\alpha, \gamma)}$.}  In any
$\mu_A$-SYT with column-$2$ entries $(\alpha, \gamma)$ on $n-2$
letters, the column-$1$ entries fill cells $(1,1), (2,1), \ldots,
(n-4,1)$ with the elements of $\{1, \ldots, n-2\} \setminus
\{\alpha, \gamma\}$ in increasing order. Consequently:
\begin{itemize}[leftmargin=2em,topsep=0.3em,itemsep=0.2em]
\item If $\gamma = n-2$: the cell $(2, 2) = n-2$ in $T'$. The
maximum column-$1$ entry is $\alpha-1$ if $\alpha = n-3$, else
$n-3$.
\item If $\gamma < n-2$: the cell $(2, 2) = \gamma$. The
maximum column-$1$ entry is $n-2$ (at $(n-4, 1)$).
\end{itemize}
For the four $(\alpha, \gamma) \in S_A$:
\begin{center}
\begin{tabular}{ccccc}
\toprule
$(\alpha, \gamma)$ & $\mathrm{pos}_{T'}(n-2)$ & $\mathrm{pos}_{T'}(n-3)$ & $\mathrm{pos}(n-1)$ in $T^{A,B}$ & $\mathrm{pos}(n)$ in $T^{A,B}$ \\
\midrule
$(n{-}5, n{-}4)$ & $(n{-}4,1)$ & $(n{-}5,1)$ & $(n{-}3, 1)$ & $(1, 3)$ \\
$(n{-}5, n{-}3)$ & $(n{-}4,1)$ & $(2, 2)$ (= $\gamma$) & $(n{-}3, 1)$ & $(1, 3)$ \\
$(n{-}4, n{-}2)$ & $(2, 2)$    & $(n{-}4, 1)$ & $(n{-}3, 1)$ & $(1, 3)$ \\
$(n{-}3, n{-}2)$ & $(2, 2)$    & $(1, 2)$ (= $\alpha$) & $(n{-}3, 1)$ & $(1, 3)$ \\
\bottomrule
\end{tabular}
\end{center}

We compute the action of $(T_{n-3}{+}1)(T_{n-2}{+}1)$ on each
$v_{T^{A,B}_{(\alpha, \gamma)}}$:

\paragraph{(A1) $(\alpha, \gamma) = (n{-}5, n{-}4)$.}
$\mathrm{pos}(n-2) = (n-4, 1)$, $\mathrm{pos}(n-1) = (n-3, 1)$ are
in the same column, adjacent. Thus $T_{n-2} v = -v$ and
\[
   (T_{n-2}{+}1) v_{T^{A,B}_{(n-5,n-4)}} = 0.
\]
Contribution to $\pi^{(c_1)} F_A$: \textbf{zero.}

\paragraph{(A2) $(\alpha, \gamma) = (n{-}5, n{-}3)$.}
Same analysis as (A1): $\mathrm{pos}(n-2) = (n-4, 1)$,
$\mathrm{pos}(n-1) = (n-3, 1)$, same column adjacent.
\[
   (T_{n-2}{+}1) v_{T^{A,B}_{(n-5,n-3)}} = 0.
\]
Contribution: \textbf{zero.}

\paragraph{(A3) $(\alpha, \gamma) = (n{-}4, n{-}2)$.}
$\mathrm{pos}(n-2) = (2, 2)$, $\mathrm{pos}(n-1) = (n-3, 1)$:
different row and column. $T_{n-2}$ acts as a $2$-block at axial
distance $d_1 = c(n-3, 1) - c(2, 2) = (4-n) - 0 = 4-n$.
For $n \ge 8$, $|d_1| \ge 4$.

By~\eqref{eq:hoefsmit}, $(T_{n-2}{+}1) v_{T^{A,B}_{(n-4,n-2)}} =
(a_{d_1}{+}1) v_{T^{A,B}_{(n-4,n-2)}} + b_{d_1} v_{\sigma_2}$,
where $\sigma_2 := s_{n-2} T^{A,B}_{(n-4, n-2)}$, obtained by
swapping the letters $n-2$ and $n-1$. Explicitly,
$\sigma_2$ has $\mathrm{pos}(n-1) = (2, 2)$,
$\mathrm{pos}(n-2) = (n-3, 1)$. (Valid SYT: row 1 reads
$(1,1) < n-4 < n$; column 1 reads $1 < 2 < \cdots < n-3 < n-2$
where the swap moved $n-2$ from $(2,2)$ to $(n-3, 1)$ and shifted
$n-3$ accordingly; column 2 reads $n-4 < n-1$.)

Now apply $T_{n-3}$:
\begin{itemize}[leftmargin=2em,topsep=0.3em,itemsep=0.2em]
\item On $v_{T^{A,B}_{(n-4,n-2)}}$: $\mathrm{pos}(n-3) = (n-4, 1)$,
$\mathrm{pos}(n-2) = (2, 2)$. $2$-block at axial distance
$d_2 = c(2, 2) - c(n-4, 1) = 0 - (5-n) = n-5$. For $n \ge 8$,
$|d_2| \ge 3$. Hence
\[
   (T_{n-3}{+}1) v_{T^{A,B}_{(n-4,n-2)}}
   = (a_{d_2}{+}1) v_{T^{A,B}_{(n-4,n-2)}} + b_{d_2} v_{\sigma_3},
\]
where $\sigma_3 := s_{n-3} T^{A,B}_{(n-4, n-2)}$:
$\mathrm{pos}(n-3) = (2, 2)$, $\mathrm{pos}(n-2) = (n-4, 1)$.
\item On $v_{\sigma_2}$: $\mathrm{pos}(n-3) = (n-4, 1)$,
$\mathrm{pos}(n-2) = (n-3, 1)$. Same column, adjacent.
$T_{n-3} v_{\sigma_2} = -v_{\sigma_2}$, so
$(T_{n-3}{+}1) v_{\sigma_2} = 0$.
\end{itemize}
Combining,
\[
   (T_{n-3}{+}1)(T_{n-2}{+}1) v_{T^{A,B}_{(n-4,n-2)}}
   = (a_{d_1}{+}1)(a_{d_2}{+}1) v_{T^{A,B}_{(n-4,n-2)}}
     + (a_{d_1}{+}1) b_{d_2} v_{\sigma_3} + 0.
\]
Both coefficients are in $\mathbb{Q}(q)^\times$.

\paragraph{(A4) $(\alpha, \gamma) = (n{-}3, n{-}2)$.}
$\mathrm{pos}(n-3) = (1, 2) = \alpha$, $\mathrm{pos}(n-2) = (2, 2) = \gamma$:
same column, adjacent. Thus $T_{n-3} v_{T^{A,B}_{(n-3,n-2)}} =
-v$, so applying $(T_{n-3}{+}1)$ to the ``identity-piece'' of the
$T_{n-2}$ $2$-block yields zero.

$\mathrm{pos}(n-2) = (2, 2)$ and $\mathrm{pos}(n-1) = (n-3, 1)$:
$2$-block at axial distance $d_1 = 4-n$, same as in (A3).
$(T_{n-2}{+}1) v_{T^{A,B}_{(n-3,n-2)}} = (a_{d_1}{+}1) v_{T^{A,B}_{(n-3,n-2)}}
+ b_{d_1} v_{\sigma_2'}$, where
$\sigma_2' := s_{n-2} T^{A,B}_{(n-3,n-2)}$:
$\mathrm{pos}(n-1) = (2, 2)$, $\mathrm{pos}(n-2) = (n-3, 1)$.

Apply $T_{n-3}$:
\begin{itemize}[leftmargin=2em,topsep=0.3em,itemsep=0.2em]
\item On $v_{T^{A,B}_{(n-3,n-2)}}$: same-column-$2$ pair $(1,2),(2,2)$,
so $(T_{n-3}{+}1) v = 0$.
\item On $v_{\sigma_2'}$: $\mathrm{pos}(n-3) = (1, 2)$,
$\mathrm{pos}(n-2) = (n-3, 1)$: different row/column.
$2$-block at axial distance
$d_3 = c(n-3, 1) - c(1, 2) = (4-n) - 1 = 3-n$, with
$|d_3| \ge 5$ at $n \ge 8$. Thus
$(T_{n-3}{+}1) v_{\sigma_2'} = (a_{d_3}{+}1) v_{\sigma_2'} +
b_{d_3} v_{\sigma_3'}$, where
$\sigma_3' := s_{n-3} \sigma_2'$: $\mathrm{pos}(n-3) = (n-3, 1)$,
$\mathrm{pos}(n-2) = (1, 2)$, $\mathrm{pos}(n-1) = (2, 2)$.
\end{itemize}
Combining,
\[
   (T_{n-3}{+}1)(T_{n-2}{+}1) v_{T^{A,B}_{(n-3,n-2)}}
   = (a_{d_1}{+}1)\cdot 0
     + b_{d_1}\,[(a_{d_3}{+}1) v_{\sigma_2'} + b_{d_3} v_{\sigma_3'}].
\]

\subsection{Disjointness of the four post-image supports}

The four contributing supports of $\pi^{(c_1)} F_A$ are:
\begin{itemize}[leftmargin=2em,topsep=0.3em,itemsep=0.2em]
\item From (A3): $\{T^{A,B}_{(n-4,n-2)},\ \sigma_3\}$.
\item From (A4): $\{\sigma_2',\ \sigma_3'\}$.
\end{itemize}
We tabulate their column-$2$ entries $((1,2),\ (2,2))$:
\begin{center}
\begin{tabular}{lccc}
\toprule
SYT & $(1, 2)$ & $(2, 2)$ \\
\midrule
$T^{A,B}_{(n-4,n-2)}$ & $n-4$ & $n-2$ \\
$\sigma_3 = s_{n-3} T^{A,B}_{(n-4,n-2)}$ & $n-4$ & $n-3$ \\
$\sigma_2' = s_{n-2} T^{A,B}_{(n-3,n-2)}$ & $n-3$ & $n-1$ \\
$\sigma_3' = s_{n-3} \sigma_2'$ & $n-2$ & $n-1$ \\
\bottomrule
\end{tabular}
\end{center}
All four $((1,2), (2,2))$ values are distinct, so the four SYTs
are distinct.

\begin{proposition}[Lemma A]\label{prop:LemmaA}
$\pi^{(c_1)} F_A \ne 0$. Explicitly,
\[
   \pi^{(c_1)} F_A
   = \tau_A\,\big[
        A_{(n-4,n-2)} (a_{d_1}{+}1)(a_{d_2}{+}1)\, v_{T^{A,B}_{(n-4,n-2)}}
        + A_{(n-4,n-2)} (a_{d_1}{+}1)\, b_{d_2}\, v_{\sigma_3}
\]
\[
        + A_{(n-3,n-2)}\, b_{d_1}(a_{d_3}{+}1)\, v_{\sigma_2'}
        + A_{(n-3,n-2)}\, b_{d_1}\, b_{d_3}\, v_{\sigma_3'}
     \big],
\]
with all four coefficients in $\mathbb{Q}(q)^\times$.
\end{proposition}

\begin{proof}
The four contributing SYTs are distinct
(table above) and each appears in $\pi^{(c_1)} F_A$ with coefficient
the product of nonzero factors $\tau_A$, the relevant
$A_{(\alpha,\gamma)}$ (nonzero by Theorem~\ref{thm:uA}), and
Hoefsmit coefficients $a_{d_i}{+}1$ or $b_{d_i}$ (nonzero
by~\eqref{eq:hoefsmit} since $|d_i| \ge 3$ in all cases at
$n \ge 8$). Hence $\pi^{(c_1)} F_A$ has four nonzero seminormal
coordinates.
\end{proof}

\section{$\pi^{(c_2)} F_B \ne 0$ (Lemma B)}\label{sec:LemmaB}

The analysis for the $B$-piece is parallel, with the key difference
that the $\mu_B$-SYT in question is a hook $(3, 1^{n-5})$, indexed
by the pair $(a, b)$ at cells $(1, 2), (1, 3)$. We write
$T^{B,B}_{(a, b)}$ for the SYT of $\lambda$ obtained from the
$\mu_B$-SYT $T_{(a, b)}'$ by $\Phi_B$ in the
``$n$ at $c_2$'' direction: $n$ at $(2, 2)$, $n-1$ at $(n-3, 1)$.

\paragraph{Positions in $T^{B,B}_{(a, b)}$.}
In $\mu_B = (3, 1^{n-5})$, the SYT $T_{(a, b)}'$ has $1$ at
$(1,1)$, $a$ at $(1, 2)$, $b$ at $(1, 3)$, and column 1
entries $\{1, \ldots, n-2\} \setminus \{a, b\}$ in increasing order
at $(1, 1), (2, 1), \ldots, (n-4, 1)$. Hence:
\begin{itemize}[leftmargin=2em,topsep=0.3em,itemsep=0.2em]
\item If $b = n-2$: cell $(1, 3) = n-2$, and the maximum
column-$1$ entry is $n-3$ if $a < n-3$, else $n-4$.
\item If $b < n-2$: cell $(1, 3) = b$, and the maximum
column-$1$ entry is $n-2$ at $(n-4, 1)$.
\end{itemize}
For the five $(a, b) \in S_B$:

\begin{center}
\begin{tabular}{ccccc}
\toprule
$(a, b)$ & $\mathrm{pos}(n-2)$ & $\mathrm{pos}(n-3)$ & $\mathrm{pos}(n-1)$ & $\mathrm{pos}(n)$ \\
\midrule
$(n{-}5, n{-}4)$ & $(n{-}4, 1)$ & $(n{-}5, 1)$ & $(n{-}3, 1)$ & $(2, 2)$ \\
$(n{-}5, n{-}3)$ & $(n{-}4, 1)$ & $(1, 3)$ (= $b$) & $(n{-}3, 1)$ & $(2, 2)$ \\
$(n{-}4, n{-}3)$ & $(n{-}4, 1)$ & $(1, 3)$ (= $b$) & $(n{-}3, 1)$ & $(2, 2)$ \\
$(n{-}4, n{-}2)$ & $(1, 3)$ (= $b$) & $(n{-}4, 1)$ & $(n{-}3, 1)$ & $(2, 2)$ \\
$(n{-}3, n{-}2)$ & $(1, 3)$ (= $b$) & $(1, 2)$ (= $a$) & $(n{-}3, 1)$ & $(2, 2)$ \\
\bottomrule
\end{tabular}
\end{center}

\paragraph{(B1, B2, B3): $b \le n-3$.}
For $(a, b) \in \{(n{-}5, n{-}4), (n{-}5, n{-}3), (n{-}4, n{-}3)\}$,
the position of $n-2$ is $(n-4, 1)$, same column as
$\mathrm{pos}(n-1) = (n-3, 1)$, adjacent.
$(T_{n-2}{+}1) v_{T^{B,B}_{(a, b)}} = 0$. Contribution: \textbf{zero.}

\paragraph{(B4) $(a, b) = (n{-}4, n{-}2)$.}
$\mathrm{pos}(n-2) = (1, 3)$, $\mathrm{pos}(n-1) = (n-3, 1)$:
$2$-block at axial distance
$d_4 = c(n-3, 1) - c(1, 3) = (4-n) - 2 = 2-n$.
For $n \ge 8$, $|d_4| \ge 6$.
$(T_{n-2}{+}1) v_{T^{B,B}_{(n-4,n-2)}} = (a_{d_4}{+}1) v_{T^{B,B}_{(n-4,n-2)}}
+ b_{d_4} v_{\rho_2}$, where
$\rho_2 := s_{n-2} T^{B,B}_{(n-4, n-2)}$: $\mathrm{pos}(n-1) = (1, 3)$,
$\mathrm{pos}(n-2) = (n-3, 1)$.

Apply $T_{n-3}$:
\begin{itemize}[leftmargin=2em,topsep=0.3em,itemsep=0.2em]
\item On $v_{T^{B,B}_{(n-4,n-2)}}$: $\mathrm{pos}(n-3) = (n-4, 1)$,
$\mathrm{pos}(n-2) = (1, 3)$. $2$-block at axial distance
$d_5 = c(1, 3) - c(n-4, 1) = 2 - (5-n) = n-3$, with $|d_5| \ge 5$
at $n \ge 8$.
$(T_{n-3}{+}1) v_{T^{B,B}_{(n-4,n-2)}} = (a_{d_5}{+}1) v_{T^{B,B}_{(n-4,n-2)}}
+ b_{d_5} v_{\rho_3}$,
$\rho_3 := s_{n-3} T^{B,B}_{(n-4, n-2)}$: $\mathrm{pos}(n-3) = (1, 3)$,
$\mathrm{pos}(n-2) = (n-4, 1)$.
\item On $v_{\rho_2}$: $\mathrm{pos}(n-3) = (n-4, 1)$,
$\mathrm{pos}(n-2) = (n-3, 1)$. Same column, adjacent.
$(T_{n-3}{+}1) v_{\rho_2} = 0$.
\end{itemize}
Contribution from (B4): $(a_{d_4}{+}1)(a_{d_5}{+}1) v_{T^{B,B}_{(n-4,n-2)}}
+ (a_{d_4}{+}1) b_{d_5} v_{\rho_3}$.

\paragraph{(B5) $(a, b) = (n{-}3, n{-}2)$.}
$\mathrm{pos}(n-3) = (1, 2) = a$, $\mathrm{pos}(n-2) = (1, 3) = b$:
same row, adjacent. Thus $T_{n-3} v_{T^{B,B}_{(n-3, n-2)}} = q v$
and $(T_{n-3}{+}1) v = (q+1) v_{T^{B,B}_{(n-3, n-2)}}$.

$\mathrm{pos}(n-2) = (1, 3)$, $\mathrm{pos}(n-1) = (n-3, 1)$:
$2$-block at axial distance $d_4 = 2-n$ (same as B4).
$(T_{n-2}{+}1) v_{T^{B,B}_{(n-3, n-2)}} = (a_{d_4}{+}1) v_{T^{B,B}_{(n-3, n-2)}}
+ b_{d_4} v_{\rho_2'}$, where
$\rho_2' := s_{n-2} T^{B,B}_{(n-3, n-2)}$: $\mathrm{pos}(n-1) = (1, 3)$,
$\mathrm{pos}(n-2) = (n-3, 1)$.

Apply $T_{n-3}$:
\begin{itemize}[leftmargin=2em,topsep=0.3em,itemsep=0.2em]
\item On $v_{T^{B,B}_{(n-3, n-2)}}$: same-row pair, so
$(T_{n-3}{+}1) v = (q+1) v$.
\item On $v_{\rho_2'}$: $\mathrm{pos}(n-3) = (1, 2)$,
$\mathrm{pos}(n-2) = (n-3, 1)$. $2$-block at axial distance
$d_3 = 3-n$ (same as the corresponding $A$-case).
$(T_{n-3}{+}1) v_{\rho_2'} = (a_{d_3}{+}1) v_{\rho_2'}
+ b_{d_3} v_{\rho_3'}$,
$\rho_3' := s_{n-3} \rho_2'$: $\mathrm{pos}(n-3) = (n-3, 1)$,
$\mathrm{pos}(n-2) = (1, 2)$, $\mathrm{pos}(n-1) = (1, 3)$.
\end{itemize}
Contribution from (B5): $(q+1)(a_{d_4}{+}1) v_{T^{B,B}_{(n-3, n-2)}}
+ b_{d_4}\,[(a_{d_3}{+}1) v_{\rho_2'} + b_{d_3} v_{\rho_3'}]$.

\subsection{Disjointness of the five post-image supports}

The five SYTs are tabulated by their values at $(1, 2)$ and $(1, 3)$:
\begin{center}
\begin{tabular}{lccc}
\toprule
SYT & $(1, 2)$ & $(1, 3)$ \\
\midrule
$T^{B,B}_{(n-4, n-2)}$ & $n-4$ & $n-2$ \\
$\rho_3 = s_{n-3} T^{B,B}_{(n-4, n-2)}$ & $n-4$ & $n-3$ \\
$T^{B,B}_{(n-3, n-2)}$ & $n-3$ & $n-2$ \\
$\rho_2' = s_{n-2} T^{B,B}_{(n-3, n-2)}$ & $n-3$ & $n-1$ \\
$\rho_3' = s_{n-3} \rho_2'$ & $n-2$ & $n-1$ \\
\bottomrule
\end{tabular}
\end{center}
All five $((1,2), (1,3))$ values are distinct, so the five SYTs
are distinct.

\begin{proposition}[Lemma B]\label{prop:LemmaB}
$\pi^{(c_2)} F_B \ne 0$. The seminormal expansion of
$\pi^{(c_2)} F_B$ has five distinct supports
$\{T^{B,B}_{(n-4, n-2)}, \rho_3, T^{B,B}_{(n-3, n-2)}, \rho_2',
\rho_3'\}$, each with coefficient in $\mathbb{Q}(q)^\times$.
\end{proposition}

\begin{proof}
Each contribution above involves a product of nonzero factors
$\tau_B$, the relevant $B_{(a, b)}$ (nonzero by Theorem~\ref{thm:uB}),
$(q+1)$, and Hoefsmit coefficients $a_{d_i}{+}1$ or $b_{d_i}$
(nonzero by~\eqref{eq:hoefsmit} since $|d_i| \ge 5$ at $n \ge 8$).
The five SYTs are distinct (table above), so the five coordinates
in $\pi^{(c_2)} F_B$ are independent and all nonzero.
\end{proof}

\section{Main theorem}\label{sec:main}

\begin{proof}[Proof of Theorem~\ref{thm:main}]
By May-20 Theorem~6.4 (the upper bound),
$\dim B_{n-2}^{(\lambda)} V_\lambda \le 2$, and the reduction
formula (Theorem~\ref{thm:reduction}) shows
$B_{n-2}^{(\lambda)} V_\lambda$ is spanned by $\{F_A, F_B\}$.

For the lower bound, suppose $c_A F_A + c_B F_B = 0$ for
$c_A, c_B \in \mathbb{Q}(q)$. Apply $\pi^{(c_1)}$:
$c_A \pi^{(c_1)} F_A + c_B \pi^{(c_1)} F_B = 0$.
By Corollary~\ref{cor:cross}, $\pi^{(c_1)} F_B = 0$, so
$c_A \pi^{(c_1)} F_A = 0$. By Proposition~\ref{prop:LemmaA},
$\pi^{(c_1)} F_A \ne 0$, hence $c_A = 0$.

Symmetrically, applying $\pi^{(c_2)}$ and using
Proposition~\ref{prop:LemmaB} gives $c_B = 0$. Therefore
$\{F_A, F_B\}$ is linearly independent and
$\dim B_{n-2}^{(\lambda)} V_\lambda = 2$.
\end{proof}

\section{Computational verification}\label{sec:verify}

All checks at $q = 7$ in $\mathbb{Q}$. Scripts at
\texttt{\textasciitilde/projects/scratch/2026-05-12-dim-2-321/}.

\paragraph{Direct dim and projection-witness check at $n \in \{8, 9, 10\}$.}
For each $n$:
\begin{itemize}[leftmargin=2em,topsep=0.3em,itemsep=0.2em]
\item Compute $F_A, F_B \in V_\lambda$ via the three-branch
reduction formula.
\item Verify $\dim \mathrm{span}(F_A, F_B) = 2 = \dim B_{n-2}^{(\lambda)} V_\lambda$
(both via direct $R'_{n-2} R'_{n-1} R'_n$ on $V_\lambda$).
\item Verify $\mathrm{supp}(F_A)$ contains only SYTs with
$\mathrm{pos}(n) \in \{(1, 3), (n-3, 1)\}$
(0-indexed: $(0, 2)$ or $(n-4, 0)$).
\item Verify $\mathrm{supp}(F_B)$ contains only SYTs with
$\mathrm{pos}(n) \in \{(2, 2), (n-3, 1)\}$.
\item Verify $\pi^{(c_1)} F_A$ has exactly $4$ nonzero coordinates,
on SYTs with $((1,2), (2,2)) \in \{(n{-}4, n{-}2), (n{-}4, n{-}3),
(n{-}3, n{-}1), (n{-}2, n{-}1)\}$ --- matching the predictions in
Proposition~\ref{prop:LemmaA}.
\item Verify $\pi^{(c_2)} F_B$ has exactly $5$ nonzero coordinates,
on SYTs with $((1,2), (1,3)) \in \{(n{-}4, n{-}2), (n{-}4, n{-}3),
(n{-}3, n{-}2), (n{-}3, n{-}1), (n{-}2, n{-}1)\}$ --- matching
Proposition~\ref{prop:LemmaB}.
\end{itemize}
All checks pass uniformly. Histograms: $|\mathrm{supp}(F_A)|$
splits as $11$ at $(n-3, 1)$ / $4$ at $(1, 3)$;
$|\mathrm{supp}(F_B)|$ splits as $13$ at $(n-3, 1)$ / $5$ at $(2, 2)$.

\section{Discussion}\label{sec:discussion}

\paragraph{The structural picture.}
This is the first proof of a rank-corner formula at $r \ge 2$. The
position-of-$n$ invariant is the conceptual core: under the
staircase $B_{\ell+1}^{(\lambda)} = R'_{\ell+1}\cdots R'_n$, the
last operator $R'_n$ uses $T_{n-1}$ which moves $n$, but
$R'_{\ell+1}, \ldots, R'_{n-1}$ all use only $T_j$ for
$j \le n-2$ and so preserve the position of $n$. Hence the
position-$n$ profile of the image after $R'_n$ is preserved through
the rest of the staircase. For $\lambda$ with $r$ length-preserving
corners, $R'_n V_\lambda = E_{n-1}^+$ decomposes by $T_{n-1}$
$+q$-eigenvectors at the various corner pairs $\{c_i, c_k\}$
involving the column-$1$ bottom $c_k$, each carrying a different
position of $n$ at $c_i \ne c_j$ for different $i$. The
position-$n$ projections separate the contributions.

\paragraph{Conjecture: rank-corner formula in general.}
We extract from the proof the following:

\begin{conjecture}[Rank-corner formula]\label{conj:rank-corner}
For any rank-zero $\lambda \vdash n$ that is not the sign
representation, with $r$ length-preserving corners and one
column-$1$ bottom corner,
\[
   \dim B_{\ell+1}^{(\lambda)} V_\lambda = r.
\]
\end{conjecture}

The structural reason: the May-19 abstract Phase A theorem +
multi-branch reduction yields $\Phi_X$ pieces for each
length-preserving corner $c_X$ paired with the column-$1$ bottom
$c_k$. Each non-trivial pair contributes a $1$-dimensional image
in $E_1^-$; pairs of length-preserving corners with each other
contribute $0$ via $j$-formula leakage (May-20 Lemma~4.1). The
position-of-$n$ invariant separates the $r$ contributions
\emph{provided} the placement cells $\{c_1, \ldots, c_{r-1}\}$ are
pairwise distinct (which they always are by definition of
corners). The non-vanishing of each piece is a finite Hoefsmit
case analysis analogous to Sections~\ref{sec:LemmaA},
\ref{sec:LemmaB} of this paper.

\paragraph{Next steps.}
\begin{enumerate}[leftmargin=2em,topsep=0.3em,itemsep=0.2em,label=(\arabic*)]
\item Prove Conjecture~\ref{conj:rank-corner} for $(4, 2, 1^{n-6})$
and $(5, 2, 1^{n-7})$ ($r = 2$) as further evidence; the technique
in this paper should adapt verbatim (different axial distances but
same case structure).
\item Handle $r = 3$: the smallest cases are $(4, 3, 2, 1^{n-9})$
and $(3, 3, 2, 1^{n-8})$ (corners at row 1, 2, 3, and column-$1$
bottom). The position-of-$n$ argument generalizes; the case
analysis grows but stays finite.
\item Connect to the underlying combinatorial object: the
length-preserving corners of $\lambda$ index ``cell-pair states''
in $E_{n-1}^+$, and the rank-corner formula counts how many of
these survive both the $j$-formula collapse and the $\mathrm{im}\subseteq E_1^-$
constraint. A bijective description in terms of partition data
would close the rank-zero programme.
\end{enumerate}

\paragraph{Status.} Theorem~\ref{thm:main} closes the dimension gap
flagged in May-20 \S 7 Remark~6.5(1). Combined with May-20
Theorem~6.4 and the May-19 abstract Phase A,
$B_{\ell+1}^{(3, 2, 1^{n-5})} V_{(3, 2, 1^{n-5})}$ is now fully
characterized: $E_1^-$ containment + exact dimension $2$.

\end{document}
